\section{The Argument: Consent as a Sorting Institution}
\subsection{The institution and its two publics}
American local governments finance capital with borrowed money, and borrowing requires authorisation. The rules of authorisation, written into state constitutions and statutes, make two separable choices. The first is scope: which debt must be put to the voters at all. Every state draws this line through the forms of finance, placing general obligation debt, backed by the taxing power, inside the requirement, and leaving revenue obligations, leases, and the debt of statutory districts and authorities outside it. The second is height: for debt inside the line, how much of which electorate must agree, a simple majority, fifty-five per cent, or two thirds. Scope decides whether a gate exists; height decides how heavy it is.

A rule of this kind creates two publics, and the difference between them is the institution's distributive core. Inside the line stands a consenting public: the electorate whose refusal stops the act. The law fixes only the share of this public that must agree. Because governments choose election dates, and timing changes who votes (Berry and Gersen 2010; Anzia 2014), the coalition actually required, the threshold multiplied by turnout, is partly constructed by the governments the rule constrains. Off-cycle timing's largest compositional effect is age, with the elderly most overrepresented in low-turnout special elections (Kogan, Lavertu and Peskowitz 2018), city elections systematically overrepresent what one study calls high-propensity voters like older, white homeowners (Hajnal, Kogan and Markarian 2022), and participation venues reliably select the same public (Hajnal, Lewis and Louch 2002; Einstein, Glick and Palmer 2019). Outside the line stands a paying public the rule creates by omission: the ratepayers, tenants and users who repay the exempt forms of debt and hold a right of refusal over none of them. The contrast is not rhetorical. The same Texas majority rule that required Harris County to assemble half a million affirmative votes authorised a developer's utility district on two, and the difference is legal architecture rather than accident. After the Supreme Court struck property qualifications in general bond elections (Cipriano v. City of Houma 1969; City of Phoenix v. Kolodziejski 1970), it upheld property-weighted and even acreage-weighted voting in specialised districts (Salyer Land Co. v. Tulare Lake Basin Water Storage District 1973; Ball v. James 1981). The freeholder franchise the general channel lost survives, lawfully, in the exit.

Three features of the requirement organise the analysis. First, it is external and frozen by comparison. The rules descend from the settlements that followed the default waves of nineteenth-century development finance, when states and localities that had borrowed for canals and railroads repudiated, restructured, and then constitutionalised restraint (Grinath, Wallis and Sylla 1997; Wallis 2005; Hillhouse 1936; Dove 2014; Sbragia 1996). The governments they bind cannot amend them, and between the only two national compilations of local debt restrictions, in 1961 and 1986, roughly a tenth of comparable entries show substantive change, over decades in which the state constitutions around these articles were amended freely (Dinan 2006; Tarr 1998). Where communities impose fiscal limits on themselves, the limits are locally adopted and locally repealable; these rules persist on nobody's continuing consent. Second, the requirement prices the form of an action, not its merit. Nothing in the rules asks whether a project is worthwhile; they ask how the money will be claimed, and the scope line tracks a single economic property, chargeability, the capacity to identify, exclude and bill a good's beneficiaries. Water, power and parking can leave the line; schools and parks cannot. Third, the price of consent is paid in a currency whose value varies across communities. Which feature of an electorate gives the requirement its force is a question with two established answers, developed below, and the evidence adjudicates between them.

\subsection{The institution in operation}
What the requirement does cannot be read from referendum returns, because most of its operation occurs before any ballot is printed and after any refusal is recorded. Research on direct democracy reached the corresponding insight for initiatives a generation ago: ballot institutions work through agendas, anticipation and strategic response, not only through votes cast (Gerber 1996; Lupia and Matsusaka 2004; Matsusaka 2004; Gerber, Lupia, McCubbins and Kiewiet 2004). The account here extends that insight to the oldest fiscal referendum in American law, tracing the requirement's operation through four decisions that every borrowing government makes in sequence.

The first decision is the channel. A government with a project asks whether it can be financed outside the voted channel at all, and the answer depends on the project's chargeability and on the menu its legal form provides. The menu is radically unequal: a school district's exits are thin, largely lease and certificate structures (Rivera 2016), a city holds many, and an authority is itself an exit. Nor is the menu given from outside; American governments multiply by design, and the creation of districts and authorities has long translated private purposes into public form (Burns 1994), with fiscal consequences that a fragmented local commons makes systematic (Berry 2009; Bunch 1991). Contemporaries saw the function early: legal commentary was describing building authorities as a costly subversion of constitutional debt limits by 1962 (Yale Law Journal 1962), and the entrepreneurial city of the postwar decades was built substantially through such devices (Sbragia 1996; Hackworth 2007). When a government takes an exit, no election occurs, and the paying public repays the project.

The second decision, where the voted channel is chosen or unavoidable, is the shape of the ask. Governments propose against a reversion, and the power to decide what reaches the ballot, how large, how bundled, and when, converts even a binding referendum into an instrument its subject partly controls. This is the agenda-setting insight of Romer and Rosenthal (1978), whose school-budget evidence showed proposers exploiting exactly this power (Romer and Rosenthal 1982), and it is why the ballot's own construction moves approval (Bechard, Lang and Pearson-Merkowitz 2024). The third decision belongs to the voters, but its record is systematically selected: because governments know their communities, they mostly submit what will pass, so high passage rates signal anticipation rather than weak constraint, and the excess of narrow victories over narrow defeats is a documented signature of fiscal elections (Cellini, Ferreira and Rothstein 2010). Anticipation is imperfect, however, and some measures fail by fractions of a point. These near-misses, accidents relative to what officials intended, are what the causal design exploits.

The fourth decision follows refusal, and it has three branches: return to the voters, reroute through an exit, or abandon. Returning is cheap, the next election is rarely a year away, and the electorate that refused is not necessarily the electorate that will answer next. Rerouting is available exactly where chargeability permits. Abandonment is the residual, rational only where a project cannot carry the cost of another campaign. The claim that organises Sections 6 and 7 is that the requirement's price is paid mainly in time and in form: refusal delays, reassigns costs from taxpayers to ratepayers, and only rarely kills.

Seen whole, the sequence adds a third discipline to the two around which the study of American local government has been organised. Residents discipline governments from below by sorting across jurisdictions (Tiebout 1956); markets and interlocal competition discipline them from above through mobile capital and the tax base (Peterson 1981; Logan and Molotch 1987; Stone 1989). Here the obligation itself moves: sideways, out of the consent requirement, across instruments and jurisdictional boundaries, without any resident or firm leaving. Which projects hold that option is set by chargeability, and it determines where a century-old requirement still binds. The nearest contemporary study shows the pattern at the city level: classifying the states into five referendum regimes, from seven with no requirement to those voting on every general obligation pledge, it finds that cities under the requirements remain active issuers but shift toward revenue and limited-tax instruments, carrying aggregate yield spreads roughly 25 to 30 basis points higher (Huang, Munevar and Samuels 2026). The vote's bite inside the regulated channel is priced by the market as well: barely approved school authorisations raise the yields of a district's existing bonds within a year, an effect absent for the cities and counties that hold the exits (Yu, Chen and Robbins 2022). What that finding leaves open, the incidence across governments, purposes and publics, and the historical construction of the requirement itself, is this paper's subject.

\subsection{Four hypotheses}
Four hypotheses follow from the account, each with a stated failure condition. H1a, the agenda: where rules are heavier, the voted share of capital should be lower, and proposals should be fewer, larger, more bundled and better timed; the hypothesis is idle if the voted share does not move with the rules. Anticipation should also have a ceiling. A bar of fifty or fifty-five per cent can be absorbed by selective proposing, so passage rates should barely differ across those regimes; at two thirds, the required coalition outruns what design can assemble for many worthwhile projects, and the requirement should begin to ration rather than shape. H1b, the response: narrow refusal should produce delay, re-submission and substitution rather than abandonment, with substitution concentrated where exits exist; the hypothesis is refuted if narrowly refused governments simply stop building. H2, composition: where requirements are heavier, capital formation should tilt away from non-chargeable purposes. This prediction turns a standing result inside out. The fiscal-institutions literature, from state budget rules (Poterba 1994; Alt and Lowry 1994; Bohn and Inman 1996; Primo 2007) to debt restrictions and their evasions (Kiewiet and Szakaly 1996; Mullins and Wallin 2004; Carr 2018), has documented substitution around limits and named it circumvention, leakage, or creative compliance (Milesi-Ferretti 2004). Treated instead as the institution's normal operation, the substitution is the mechanism of its distributive work, and, Section 2.4 argues, of its survival.

H3 concerns the electorate, and two established accounts disagree in a way the data can settle. On an assembly-cost reading, the requirement should bind hardest where communities are divided, because social heterogeneity raises the price of every coalition; the canonical finding that ethnically fragmented places provide fewer shared public goods supplies the premise (Alesina, Baqir and Easterly 1999; Alesina, Glaeser and Sacerdote 2001). On a transmission reading, a vote binds only where the electorate is durable enough that today's answer is also tomorrow's, so the requirement should be most consequential where the public convened is stable, propertied and homogeneous: where refusals stay refused and approvals license years of borrowing. The premise here is the homevoter, the owner-occupier whose undiversified asset makes local fiscal choices personal and whose attention makes them enforceable (Fischel 2001; Oliver 2001). The two readings disagree about the sign on social heterogeneity, which is what makes the adjudication informative, and neither originates with this paper.

\subsection{Why the rules persist}
If the requirement binds anyone, why has almost no one amended it in a century? The literature on institutional durability explains persistence through increasing returns and sequencing (Pierson 2004) and change through its quieter modes: displacement, layering, drift, conversion (Streeck and Thelen 2005; Hacker 2004; Mahoney and Thelen 2010; Orren and Skowronek 2004). The public-authority statutes are layering; the judicial narrowing of ``debt'' is conversion; the unadjusted nominal limits are drift. The mechanism proposed here supplies what those categories describe but do not explain: the feedback that protects the core rule while everything moves around it. Amending a constitution requires a statewide coalition, and the natural members are the interests the gate blocks. But any blocked interest whose purpose is chargeable can exit for the price of a bond resolution, and each exit removes a member from the reform coalition. The rule persists through selective escape: it leaks exactly enough to drain the pressure that would amend it, while the population still exposed to it becomes, over time, the set of purposes and places with no alternative. The prediction is sharp, and Section 9 reads the modern record of threshold politics against it: campaigns to change a bar should arise from, and succeed only in, the sector the scope line offers no exit. The mechanism also carries a general amendment to the theory of veto points: formally identical vetoes are not equivalent institutions, because a veto's force is set by the outside options that surround it (cf. Tsebelis 2002).

\subsection{The stakes}
The distributive stakes follow from the structure. Because scope tracks chargeability and exit capacity is chargeability's twin, the weight of the regime falls on the unpriced goods of the local state, education first among them, goods whose blocked investment is worth more than it costs, since marginal homebuyers value each dollar of narrowly approved school spending at 1.50 dollars or more (Cellini, Ferreira and Rothstein 2010), while priced goods, including services for poor users delivered in enterprise form, are supplied without consent and financed by flat charges. The perimeter this creates is the local instance of a state built out of sight (Balogh 2009; Mettler 2011), and its financing has never been distributively innocent: municipal debt has carried the geography of American racial inequality in its prices and its access (Jenkins 2021), property-tax constitutionalism carries older exclusions still (Einhorn 2006), and several of the fiscal constitutions studied here were written in conventions whose disenfranchising purposes were explicit [V: Kousser; Alabama 1901]. The supermajority variants sharpen the question. In the constitutional tradition, higher thresholds trade decision costs for protection against the external costs majorities impose (Buchanan and Tullock 1962); critics answer that supermajority rules entrench whatever the status quo distributes and can hand minorities government by inaction (McGann 2004; Schwartzberg 2014). That dispute is usually conducted normatively. Here it becomes a measurement: Section 8 identifies the minority the bar in fact protects, propertied, older, homogeneous electorates, with officeholder partisanship null here as in municipal capital composition generally (Hirt 2026), and finds the blocked majorities in poorer districts, communities whose local institutions have long been instruments of composition rather than mirrors of it (Trounstine 2018). And the collision runs through the constitutions themselves: the documents that command educational provision, enforceably so in the school-finance cases (Zackin 2013; Rose v. Council for Better Education 1989 [V]), are the same documents that gate the capital the command requires, administered through the most numerous democracies in America (Berkman and Plutzer 2005).

Three simplifications are deliberate, and the evidence licenses each. The account contains no bargaining between governments and voters over project design after refusal, because returning measures ask for the original amount and keep their purposes. It contains no mechanism by which refusal reveals information that kills projects, because refused projects are financed anyway a third of the time and extinguished less than a fifth. And it offers no welfare judgment: the aim is to establish what the institution does, for whom, and why it has lasted, not whether it should. The empirical sections follow the sequence of operation: the landscape and the channel (Section 5), the vote's effect at the threshold (Section 6), the response to refusal (Section 7), the adjudication between the two readings of H3 (Section 8), and the shaping of agendas and of the rules themselves (Section 9).

